\documentclass[../../main.tex]{subfiles}

\begin{document}
\begin{CJK*}{UTF8}{gbsn}

\section*{Exercise 1}

Consider split conformal prediction 
with data set $(X_1, Y_1), \cdots, (X_{n+1}, Y_{n+1}) \in \mathcal{X} \times \mathcal{Y}$ 
and the conformal risk control with given score 
$s : \mathcal{X} \times \mathcal{Y} \to \mathbb{R}$.
We write $S_i = s(X_i, Y_i)$ for all $i$ for simplicity.
Instead of what is in Chapter 10 of the book, we define,
for all $\lambda \geqslant 0$:

\begin{equation*}
    \mathcal{C}_{\lambda}(x)
    = \curly{y \in \mathcal{Y} \mid s(x,y) < \lambda }
\end{equation*}
and 

\begin{equation*}
    L(y,\mathcal{C}_{\lambda}(x))
= \chi_{ \curly{y \notin \mathcal{C}_{\lambda}(x)} }.
\end{equation*}

Show that, with this definition,
$L(y,\mathcal{C}_{\lambda}(x))$ 
is no longer always right continuous for every 
fixed $(x,y) \in \mathcal{X} \times \mathcal{Y}$.
Next, given $\alpha \in (0,1)$, give a concrete expression for 

\begin{equation*}
    \hat{\lambda} = \inf \curly{\lambda \geqslant 0 \mid \hat{R}(\lambda) 
    \leqslant \alpha - \frac{\alpha - 1}{n}}
\end{equation*}
in terms of the data,
where 

\begin{equation*}
    \hat{R}(\lambda) = \frac{1}{n} \sum_{i=1}^n L(Y_i,\mathcal{C}_{\lambda}(X_i)).
\end{equation*}
Show that, if we assume that $(S_1, \ldots, S_{n+1})$ is exchangeable 
and that the $S_i$'s almost surely has no tie,
then 

\begin{equation*}
    P\curly{Y_{n+1}\in\mathcal{C}_{\hat{\lambda}}(X_{n+1})}
    \geqslant 1-\alpha.
\end{equation*}
Provide a counterexample to show that 
this need not hold if there are ties in $S_i$'s with positive probability,
even if they are exchangeable.

\paragraph{Solution}
Let $(x,y) \in \mathcal{X} \times \mathcal{Y}$ be given.
If $\{\lambda_m\}_{m=1}^{\infty}$ is any decreasing sequence such that
$\lambda_m > s(x,y)$ for every $m$ and $\lambda_m \to s(x,y)$ as $m \to \infty$,
then $L(y,\mathcal{C}_{\lambda_m}(x)) = 1$ for all $m$ 
but $L(y,\mathcal{C}_{s(x,y)}(x)) = 0$, proving that 
$L(y,\mathcal{C}_{\lambda}(x))$ is not always right continuous.

We move to compute $\hat{\lambda}$.
The definition of $\hat{R}(\lambda)$ and simple arithmetic rearrangement 
shows that,
for any $\lambda \geqslant 0$,

\begin{equation*}
    \hat{R}(\lambda) \leqslant \alpha - \frac{\alpha - 1}{n}
\end{equation*}
if and only if 

\begin{equation*}
    \sum_{i=1}^n \chi_{\curly{S_i < \lambda}} \geqslant (n+1)(1-\alpha).
\end{equation*}
Thus, it is obvious that 

\begin{equation*}
    \hat{\lambda} = S_{\bracket{\lceil(n+1)(1-\alpha)\rceil}}.
\end{equation*}
Here
we have used order statistics $S_{(1)} \leqslant \cdots \leqslant S_{(n)}$
for the first $n$ values of $S_1, \ldots, S_n$.
Consequently:

\begin{equation*}
    C_{\hat{\lambda}}(X_{n+1})
    = \curly{y \in \mathcal{Y} \mid s(X_{n+1}, y) < 
    S_{\bracket{\lceil(n+1)(1-\alpha)\rceil}}}.
\end{equation*}

To prove coverage, let $\alpha \in (0,1)$ be given.
Assume first that
that $(S_1, \ldots, S_{n+1})$ is exchangeable and 
there is almost surely no tie in the $S_i$'s.
Then

\begin{equation*}
    P\curly{Y_{n+1} \in C_{\hat{\lambda}}(X_{n+1})}
    = P\curly{S_{n+1} < S_{\bracket{\lceil(n+1)(1-\alpha)\rceil}}}
    = \frac{\lceil(n+1)(1-\alpha)\rceil}{n + 1} \geqslant 1 - \alpha,
\end{equation*}
where the second equality uses both assumptions:
when $S_{(1)} < \cdots < S_{(n)}$, inserting a new exchangeable 
point $S_{n+1}$ to get exactly right before $\lceil(n+1)(1-\alpha)\rceil$th place 
has the probability $\frac{\lceil(n+1)(1-\alpha)\rceil}{n + 1}$.
If we remove the assumption that there are no ties,
then everything can go wrong. For example,
just setting $s \equiv 0$.
Then $\hat{\lambda} = 0$ and $C_{\hat{\lambda}}(x)= \emptyset$ for all $x$,
clearly not yielding any good coverage probability.
////

\end{CJK*}
\end{document}
